\documentclass[a4paper]{umbruch-position}
\usepackage[utf8]{inputenc}
\usepackage[T1]{fontenc}
\usepackage[ngerman]{babel}
\usepackage{tikz}
\usepackage{amssymb}
\usepackage{eurosym}
\usepackage{graphicx}
\usepackage{xcolor}
\definecolor{umbruch-dark}{HTML}{2A2A2A}
\definecolor{umbruch-teal}{HTML}{008080}
\definecolor{umbruch-red}{HTML}{B22222}
\usetikzlibrary{shapes.geometric, arrows.meta, positioning, fit, calc}

\positionsnummer{Basiert auf der Studie ``Systemtheoretische Aufarbeitung der Regressangst'' (ST-001)}
\versionsnummer{0.2}
\title{Mein Arzt hat Angst}
\untertitel{Der Patienten-Kompass: Wie Sie Behandlungsblockaden durch\\Regressangst aufl\"osen und Ihren Arzt zum Verb\"undeten machen}

\begin{document}

\maketitle

\begin{infobox}
\textbf{WICHTIGER HINWEIS:}\\
\textit{Dieser Leitfaden dient der allgemeinen Information und stellt keine Rechtsberatung dar. F\"ur individuelle rechtliche Beratung wenden Sie sich an einen Fachanwalt f\"ur Sozialrecht oder die Unabh\"angige Patientenberatung (UPD, 0800 011 77 22).}
\end{infobox}

\section{Das Problem in einfacher Sprache}

Ein \glqq{}Regress\grqq{} ist eine R\"uckforderung, die Ihr Arzt an die Krankenkasse zur\"uckzahlen muss, wenn diese beschlie\ss{}t, dass er \"uber das vorgegebene Budget hinaus Rezepte oder Heilmittel verordnet hat.
Das Ergebnis: Obwohl eine Behandlung medizinisch f\"ur Sie das Beste w\"are, ist sie f\"ur den Arzt mitunter existenzgef\"ahrdend. Er f\"urchtet pers\"onliche Strafzahlungen von tausenden Euro. Die \"offentlich kommunizierte Regressquote von ``unter 1\,\%'' bezieht sich auf das statistische Screening (V1) --- bei der Einzelfallpr\"ufung (V2) ist die tats\"achliche Durchsetzungsrate deutlich h\"oher, aber sie findet \"uber au\ss{}ergerichtliche Vergleiche statt und taucht in keiner Statistik auf. \textbf{Die Angst Ihres Arztes ist berechtigt} --- sie bezieht sich auf ein anderes Verfahren als die beruhigende Statistik.

\section{Warum \glqq{}1\,\% Regressrisiko\grqq{} nicht bedeutet, was Sie denken}

Wenn Sie h\"oren, dass das Regressrisiko unter 1\,\% liegt, bezieht sich diese Zahl auf das statistische Screening (Auff\"alligkeitspr\"ufung)~--- nicht auf das Verfahren, das \"Arzte tats\"achlich f\"urchten. Vier Zahlenebenen werden in der \"Offentlichkeit systematisch vermischt:

\begin{center}
{\small
\begin{tabular}{p{4cm} p{8.5cm}}
\toprule
\textbf{Was gemessen wird} & \textbf{Was das bedeutet} \\
\midrule
Auff\"alligkeitspr\"ufung eingeleitet & Statistisches Screening~--- viele Auff\"alligkeiten sind normal und gewollt; niedrige Trefferquote ist systemimmanent \\[3pt]
Einzelfallpr\"ufung eingeleitet & Konkreter Verdacht einer Kasse~--- rechtlich begr\"undet; kein Beratungsschutz. Bei V2b (Formfehler): Sofortige Sanktionslogik. Bei V2a (Unwirtschaftlichkeit): Medizinische Verteidigung und Kostendifferenz m\"oglich. \\[3pt]
Regress festgesetzt & Formale Entscheidung~--- oft im Widerspruchsverfahren anfechtbar \\[3pt]
Regress durchgesetzt & Tats\"achlich bezahlter Betrag \\
\bottomrule
\end{tabular}
}
\end{center}

\begin{center}
  \includegraphics[width=0.90\textwidth]{C:/Users/User/OneDrive/.UMBRUCH/Projekte/Regress-Melder/grafiken/vier_messebenen}
\end{center}

{\small
\begin{tabularx}{\textwidth}{p{3.5cm} X X}
\toprule
 & \textbf{Auff\"alligkeitspr\"ufung} & \textbf{Einzelfallpr\"ufung} \\
\midrule
Charakter & Such- und Filterverfahren & Gezieltes Verdachts\-verfahren \\
Beratungs\-schutz & Vorgeschrieben & Gesetzlich ausgeschlossen \\
Neg.\,Ausgang & Niedrig wahrscheinlich & Hoch wahrscheinlich \\
Schaden\-h\"ohe & Niedrig bis moderat & Moderat bis hoch \\
Wiederholbar & Eingeschr\"ankt & M\"oglich \\
\bottomrule
\end{tabularx}
}

Das Screening hat eine niedrige Trefferquote~--- das ist kein Fehler, sondern Absicht. Das Einzelfallverfahren hat --- insbesondere bei Formfehlern (V2b) --- eine \textbf{hohe} Erfolgsquote f\"ur die Kassen. Es l\"auft ohne Beratungsschutz. Bei Formfehlern direkt in die Sanktion; bei Unwirtschaftlichkeit (V2a) sind medizinische Einw\"ande m\"oglich. \textbf{Genau das f\"urchten \"Arzte.}

Ihr Arzt hat recht: Die ``unter 1\,\%'' z\"ahlen nur V1-Regresse. Bundesweit gab es 2022 aber ca.\,47.000 V2-Verfahren (BMG). In Bayern allein: 13.332 Regresse bei $\sim$24.000 \"Arzten (SpiFa 2024). \textbf{Warum trotzdem 72\,\% der Haus\"arzte betroffen sind:} Bei V2 ist das Risiko pro \textit{einzelner Verordnung} gering --- aber ein Arzt stellt pro Quartal ca.\,3.000 Verordnungen aus. \"Uber 10 Jahre sind das 120.000 Angriffspunkte. \textbf{Kleines Risiko pro Fall, gro\ss{}es Risiko pro Arzt.}

Aber auch V1 belastet: Alle markierten \"Arzte tragen die psychologischen Kosten des Screenings, auch wenn das Schutzstufenmodell (Anh\"orung, Beratung, Karenzzeit) die allermeisten entlastet.

\section{Woran Sie erkennen, dass es um Budget geht}

Wenn der Regress-Druck zuschl\"agt, h\"oren Patienten meist diese Phrasen:
\begin{itemize}
    \item \textit{\glqq{}Das ist zu teuer.\grqq{}}
    \item \textit{\glqq{}Das kann ich in diesem Quartal nicht mehr verordnen.\grqq{}}
    \item \textit{\glqq{}Wir probieren erstmal etwas G\"unstigeres, bevor wir es so kräftig ausschöpfen.\grqq{}}
    \item Oder der Arzt verweist Sie auf einen Facharzt f\"ur Rezepte, die er selbst ausstellen d\"urfte ("Verordnungsweiterleitung").
\end{itemize}

\section{Entscheidungsbaum f\"ur Patienten: \glqq{}Mein Arzt verordnet nicht\grqq{}}

\begin{infobox}
\textit{Ihre Gesundheit geht vor: Sehen Sie Ihren Arzt als Verb\"undeten gegen das System, nicht als Feind.}
\end{infobox}

\begin{center}
\tikzset{
    action/.style={rectangle, draw=umbruch-dark, fill=umbruch-teal!15, rounded corners, minimum height=1cm, align=center, text width=4cm, inner sep=0.2cm},
    decision/.style={diamond, aspect=2.0, draw=umbruch-dark, fill=white, align=center, text width=3cm, inner sep=0cm},
    warning/.style={rectangle, draw=umbruch-red, fill=umbruch-red!15, rounded corners, minimum height=1cm, align=center, text width=4.5cm, inner sep=0.2cm},
    arrow/.style={-{Stealth[scale=1.2]}, thick, umbruch-dark}
}
\begin{tikzpicture}[
    scale=0.8, transform shape,
    node distance=1.5cm and 1cm,
    font=\sffamily\small
]
\node[decision] (ablehnung) {Arzt lehnt ab!\\Frage: \glqq{}Medizinisch oder budget\"ar?\grqq{}};
\node[action, right=1.5cm of ablehnung] (med) {Medizinisch: Akzeptieren oder\\anderen Arzt befragen\\(freie Arztwahl)};
\node[decision, below=1.5cm of ablehnung] (budget) {Budget\"ar:\\"K\"onnen Sie das in\\Akte schreiben?\grqq{}};
\node[action, right=1.5cm of budget] (dokumentiert) {Dokumentiert: Gut!\\Nachweis f\"ur Arzt/Kasse\\(\it BSG B 6 KA 27/12 R)};
\node[decision, below=1.5cm of budget] (kasse) {Genehmigungs-\\pflichtig?\\(Cannabis, Reha,\\Hilfsmittel)};
\node[action, right=1.5cm of kasse] (genehmigt) {JA: Eigenantrag!\\Kasse muss in 3 Wo.\\entscheiden (Fiktion)};
\node[warning, below=1.5cm of kasse] (abgelehnt) {NEIN (Standard-\\Arzneimittel):\\Kasse verweist auf\\Arzt (\S\,29 BMV-\"A)};
\node[warning, right=1.5cm of abgelehnt] (sg) {Sozialgericht (kostenlos für GKV-Patienten!)};
\draw[arrow] (ablehnung) -- node[above] {Medizinisch} (med);
\draw[arrow] (ablehnung) -- node[left] {Budget} (budget);
\draw[arrow] (budget) -- node[above] {JA} (dokumentiert);
\draw[arrow] (budget) -- node[left] {NEIN} (kasse);
\draw[arrow] (kasse) -- node[above] {Erfolg} (genehmigt);
\draw[arrow] (kasse) -- node[left] {Abgelehnt} (abgelehnt);
\draw[arrow] (abgelehnt) -- (sg);
\end{tikzpicture}
\end{center}

\newpage
\section{Ihre Rechte als Patient}

\begin{tabularx}{\textwidth}{p{3.5cm} X}
\toprule
\textbf{Recht} & \textbf{Beschreibung \& Gesetz} \\
\midrule
Versorgungsanspruch & Sie haben Recht auf ausreichende/zweckm\"a\ss{}ige Therapie (\S\,12, \S\,27 SGB\,V). \\
Freie Arztwahl & Sie d\"urfen jederzeit einen anderen Arzt aufsuchen und eine weitere Meinung einholen. \\
Akteneinsicht & L\"uckenloser Einblick in Patientenakte (BGB \S\,630g). \\
Genehmigungsfiktion (nur bei genehmigungspflichtigen Leistungen) & Entscheidet die Kasse auf Ihren Antrag hin nicht binnen 3 Wochen (oder 5 mit MD-Gutachten), gilt die Leistung als genehmigt (\S\,13 Abs.\,3a SGB\,V). \textbf{Gilt nicht f\"ur normale Arzneimittel} --- dort verweist die Kasse auf die \"arztliche Verordnungshoheit (\S\,29 BMV-\"A). \\
\bottomrule
\end{tabularx}

\section{Zwei vergessene Wege: Was Ihr Arzt noch tun kann}

Das System hat zwei pr\"aventive Schutzinstrumente, die kaum bekannt sind. Als Patient k\"onnen Sie Ihren Arzt aktiv darauf hinweisen:

\begin{infobox}
\textbf{Weg 1: KV-Verordnungsberatung (kostenlos)}\\
Ihr Arzt kann bei der Kassen\"arztlichen Vereinigung eine kostenlose pharmakologische Beratung zu Ihrer konkreten Situation anfragen. Bei der KVBW zum Beispiel unter Telefon 0711\,7875-3663. Die Beratung ist unverbindlich, st\"arkt aber seine Dokumentation erheblich.\\[4pt]
\textbf{Was Sie sagen k\"onnen:} \textit{,,Gibt es eine M\"oglichkeit, dass Sie vorab eine Verordnungsberatung bei Ihrer KV anfragen? Das w\"are kostenlos und w\"urde Ihnen Sicherheit geben.''}
\end{infobox}

\begin{infobox}
\textbf{Weg 2: Anfrage zur Verordnungsf\"ahigkeit bei der Kasse (BSG B\,6\,KA\,27/12\,R)}\\
Ihr Arzt kann die Krankenkasse schriftlich fragen, ob eine bestimmte Verordnung erstattungsf\"ahig ist. Das BSG hat diesen Weg erlaubt. Wenn die Kasse positiv antwortet, kann sie sp\"ater keinen Regress beantragen, ohne sich selbst zu widersprechen.\\[4pt]
\textbf{WICHTIGE EINSCHR\"ANKUNG:} Dieser Weg funktioniert \textbf{nur bei unklarer Verordnungsf\"ahigkeit} --- z.\,B. Off-Label-Use, neue Wirkstoffe, Grenzf\"alle. Bei \textbf{regul\"ar verordnungsf\"ahigen Arzneimitteln} (also den meisten Standardmedikamenten) wird die Kasse antworten: \textit{,,Wir d\"urfen Verordnungen nicht vorab genehmigen''} (\S\,29 Abs.\,1 BMV-\"A). In diesem Fall ist die Anfrage vertane M\"uhe --- die Verordnungsentscheidung liegt dann vollst\"andig beim Arzt.\\[4pt]
\textbf{Wann lohnt es sich?} Wenn Ihr Arzt unsicher ist, \textit{ob} er ein Medikament \"uberhaupt verordnen darf (nicht ob er es sich leisten kann). Typisch: Cannabis, teure Biologika, Medikamente au\ss{}erhalb der zugelassenen Indikation.\\[4pt]
\textbf{Was Sie sagen k\"onnen:} \textit{,,Ist bei meinem Medikament die Verordnungsf\"ahigkeit unklar? Dann k\"onnten Sie die Kasse nach BSG-Urteil um eine Stellungnahme bitten.''}
\end{infobox}

\begin{infobox}
\textbf{Wichtig --- drei verschiedene Dinge, nicht verwechseln:}\\
(1)~\textbf{Ihr Eigenantrag bei der Kasse} (als Patient): Sie k\"onnen bei Ihrer Kasse einen Antrag auf Kosten\"ubernahme stellen. Bei genehmigungspflichtigen Leistungen (Cannabis, Reha, Hilfsmittel) muss die Kasse innerhalb von 3 Wochen entscheiden --- sonst Genehmigungsfiktion (\S\,13 Abs.\,3a SGB\,V). \textbf{Bei normalen Arzneimitteln} wird die Kasse antworten: \textit{``Die Verordnungsentscheidung liegt beim Arzt''} (\S\,29 BMV-\"A) --- hier l\"auft der Eigenantrag ins Leere.\\[2pt]
(2)~\textbf{KV-Verordnungsberatung} (Weg 1, Arzt$\to$KV): Kostenlos, st\"arkt Dokumentation. Immer m\"oglich.\\[2pt]
(3)~\textbf{KK-Anfrage Verordnungsf\"ahigkeit} (Weg 2, Arzt$\to$Kasse): Sinnvoll wenn die Verordnungsf\"ahigkeit \textit{unklar} ist (Off-Label, neue Wirkstoffe). Bei Standardmedikamenten inhaltsleer, weil die Verordnungsf\"ahigkeit ja feststeht.\\[2pt]
Weg 1 und 2 stehen nur dem Arzt offen. Ihr Eigenantrag (1) ist unabh\"angig davon.
\end{infobox}

\section{Gespr\"achsleitfaden f\"ur das Arztgespr\"ach}

Drei S\"atze, die konkret helfen k\"onnen, Blockaden abzubauen:
\begin{enumerate}
    \item \textit{\glqq{}Gibt es hier eine leitliniengerechte Behandlungs-Alternative, die wir statt-dessen anwenden k\"onnten?\grqq{}}
    \item \textit{\glqq{}K\"onnen Sie die aktuelle Aussetzung in meiner Akte bitte ausf\"uhrlich dokumentieren?\grqq{}}
    \item \textit{\glqq{}Braucht dieses Medikament eine Kassengenehmigung? Falls ja --- kann ich den Antrag selbst stellen?\grqq{}} (Nur bei genehmigungspflichtigen Leistungen wie Cannabis, Reha, Hilfsmitteln. Bei normalen Arzneimitteln verweist die Kasse auf den Arzt.)
\end{enumerate}
\textbf{Was Sie vermeiden sollten:} Vorw\"urfe, das Verlangen auf unm\"ogliche Verordnungen oder Drohungen.

\medskip
\textbf{Wirksamerer Gespr\"achsframe:} Nicht \textit{\glqq{}Bitte verschreiben Sie mir das.\grqq{}} Sondern: \textit{\glqq{}Wie k\"onnen wir Sie dabei absichern?\grqq{}} Dieser Wechsel verlagert die Frage von der Konfrontation zur gemeinsamen Probleml\"osung~--- und \"offnet den Raum f\"ur Option B und C aus dem vorigen Abschnitt.

\section{Anlaufstellen}
\begin{itemize}
    \item Unabh\"angige Patientenberatung Deutschland (UPD): \textbf{0800 011 77 22}
    \item Patientenbeauftragte der jeweiligen Krankenkasse
\end{itemize}

\section{Was Sie wissen sollten: Das System pr\"uft nur in eine Richtung}

\begin{infobox}
\textbf{Wichtig zu verstehen:} Ein Regress gegen Ihren Arzt bedeutet \textbf{nicht}, dass die Behandlung medizinisch falsch war. Das Pr\"ufsystem pr\"uft \"uberwiegend formale Regeln und Dokumentation --- nicht medizinische Qualit\"at (etwa ein Drittel der Pr\"ufungen betrifft allerdings Regeln, die auf medizinischer Evidenz basieren). Ihr Arzt kann medizinisch alles richtig gemacht haben und trotzdem zahlen m\"ussen, weil ein Formular falsch ausgef\"ullt war.
\end{infobox}

Eine wichtige strukturelle Information: Das deutsche Wirtschaftlichkeitspr\"ufsystem pr\"uft ausschlie\ss{}lich, ob \textit{zu viel} verordnet wurde. Es pr\"uft \textbf{nie}, ob \textit{zu wenig} verordnet wurde:

\begin{itemize}
    \item Wurde \textbf{zu viel} verordnet $\to$ Regress
    \item Wurde \textbf{zu wenig} verordnet $\to$ keine Konsequenz
    \item Wurden Sie \textbf{abgeschoben} statt behandelt $\to$ keine Konsequenz
    \item Wurde Ihre Therapie \textbf{verz\"ogert} aus Angst $\to$ keine Konsequenz
\end{itemize}

\textbf{Die Asymmetrie ist systemisch.} Ein Arzt, der aus Angst nichts verordnet, wird nie gepr\"uft. Ein Arzt, der seinen Patienten leitliniengerecht behandelt, riskiert Regress. Das System belohnt Unterlassung und bestraft Handeln. Diese Logik trifft Sie als Patient direkt --- sie ist nicht Ihre Schuld und nicht die Schuld Ihres Arztes. Sie ist ein strukturelles Problem.

Wenn Sie unter dieser Logik leiden, ist es wichtig zu wissen: \textbf{Sie sind nicht allein.} Eine Befragung von 770 deutschen Haus\"arzten und Orthop\"aden hat 2023 gezeigt: 37\,\% der Haus\"arzte berichten, ihr \"arztliches Handeln werde durch Regressangst ``stark beeinflusst''. 51\,\% der Haus\"arzte \"uberweisen Patienten an Spezialisten, obwohl sie selbst behandeln k\"onnten. 72\,\% haben mindestens einmal einen Regress erlebt (Ribbat et al. 2023, Gesundheitswesen 85: 111--118; V1+V2 zusammen). Ihr Arzt ist Teil eines Systems, das diese Logik seit 1989 nicht aufgel\"ost hat. Wer sich angesichts dessen besser verteidigen will, findet im Beratungsangebot der UPD und der Patientenbeauftragten konkrete Hilfe.

\section{Der Ernstfall --- Wie System-Tools Leben sch\"utzen}

\begin{infobox}
\textbf{Fallbericht: 491.000\,EUR Regress --- wegen eines Stempels}\\
\textbf{Das Problem:} Ein niedergelassener Arzt versorgte \"uber Jahre hunderte chronisch kranke Patienten. Alle Verordnungen waren medizinisch korrekt und leitliniengerecht. Doch auf den Rezepten befand sich ein Namensstempel statt einer pers\"onlichen Unterschrift.\\[4pt]
\textbf{Die Blockade:} Die Krankenkasse beantragte einen Einzelfallregress \"uber 491.000\,EUR --- f\"ur s\"amtliche Verordnungen des betroffenen Zeitraums. Nicht weil die Medizin falsch war, sondern weil die Form nicht stimmte. Das BSG best\"atigte den vollen Regress (B\,6\,KA\,9/24\,R, 27.08.2025). Das Argument des Arztes, alle Verordnungen seien medizinisch indiziert gewesen, wurde als ``irrelevant'' verworfen.\\[4pt]
\textbf{Was das f\"ur die Patienten bedeutet:} Die Patienten dieses Arztes verlieren ihren eingespielten Versorger. Die Kontinuit\"at der Behandlung ist unterbrochen --- durch ein formales Kriterium, nicht durch medizinisches Versagen.\\[4pt]
\textbf{Wie System-Tools geholfen h\"atten:}
\begin{itemize}[nosep]
  \item \textbf{Intervention 1:} Ein Formcheck-Prozess im Backoffice (Vier-Augen-Prinzip) h\"atte den Stempel bei \textit{jeder einzelnen Verordnung} vor Ausgabe erkannt.
  \item \textbf{Intervention 2:} Als Patient h\"atten Sie Ihren Arzt auf die KV-Verordnungsberatung hinweisen k\"onnen --- die h\"atte das Dokumentationsproblem aufgedeckt.
\end{itemize}
\textit{Quelle: BSG B\,6\,KA\,9/24\,R, 27.08.2025. KBV nannte das Urteil ``unverh\"altnism\"a\ss{}ig und \"uberzogen''. Verfassungsbeschwerde angek\"undigt.}
\end{infobox}

\vspace{1cm}

\begin{infobox}
\textbf{Was Sie mitnehmen k\"onnen:}\\
Niedrige Regressquoten gelten f\"ur das statistische Screening~--- nicht f\"ur das Einzelfallverfahren, das \"Arzte wirklich f\"urchten. Die Angst Ihres Arztes ist rational. Und: Es gibt Wege, ihn abzusichern~--- wenn Sie und er sie gemeinsam nutzen.
\end{infobox}

\vfill
\begin{center}
\footnotesize{\textit{Dieser Leitfaden dient der allgemeinen Information und stellt keine Rechtsberatung dar. F\"ur individuelle rechtliche Beratung wenden Sie sich an einen Fachanwalt f\"ur Sozialrecht oder die Unabh\"angige Patientenberatung (UPD, 0800 011 77 22).}}
\end{center}


\end{document}
